\section{От усиления закон към скоростта на сходимост}

За да поставим новия въпрос за скоростта на сходимост, първо припомняме означенията от усиления закон. Разглеждаме редица от независими и еднакво разпределени (н.е.р.) случайни величини $(X_i)_{i=1}^\infty$. Фактът, че са взаимно независими и еднакво разпределени, означава, че всяка една от тях има същата функция на разпределение като всяка друга, тоест функцията на разпределение $F(x)$ е една и съща за всяко $X_i$.

Нека $\mu = \E[X_1]$ бъде математическото очакване на първата случайна величина. Тъй като величините са еднакво разпределени, $\mu$ е математическото очакване на всички останали величини в редицата. Дефинираме парциалната сума $S_n$:
\[ S_n = \sum_{j=1}^n X_j \]
$S_n$ представлява нова случайна величина за всяко $n$. Редицата от случайни величини $S_1, S_2, \dots$ разбира се не е нито от независими, нито от еднакво разпределени величини, тъй като те натрупват стойностите на първите $n$ елемента.

Съгласно \crosslecture{thm:slln}{теоремата за усиления закон за големите числа от предходната лекция}, средното аритметично $\frac{S_n}{n}$ се схожда почти сигурно (п.с.) към математическото очакване $\mu$:
\[ \frac{S_n}{n} \xrightarrow[n \to \infty]{\text{п.с.}} \mu \]
Това означава, че ако извършим серия от експерименти, средното аритметично от резултатите ще се доближава неограничено до очакването на първата случайна величина с вероятност единица. Ако означим разликата (или грешката) с $E_n = \frac{S_n}{n} - \mu$, то е ясно, че $E_n$ клони към нула почти сигурно при $n \to \infty$.

Въпросът обаче е: \emph{колко бързо} се осъществява тази сходимост? Можем ли да кажем нещо допълнително за поведението на грешката $E_n = \frac{S_n}{n} - \mu$? Знаем, че клони към нула, но това е единствената информация дотук. За да отговорим на този въпрос и да изследваме поведението на тази грешка, се нуждаем от Централната гранична теорема (ЦГТ).

\section{Централна гранична теорема (ЦГТ)}

Нека формулираме теоремата, която описва прецизно поведението на тази разлика при голямо $n$.

\begin{keythm}[Централна гранична теорема]\label{thm:clt}
Нека $(X_i)_{i=1}^\infty$ са независими и еднакво разпределени (н.е.р.) случайни величини със средно $\mu = \E[X_1]$ и дисперсия $\Var[X_1] = \sigma^2 < \infty$. Тогава за
\[ Z_n := \frac{\sqrt{n}}{\sigma} E_n = \frac{\sqrt{n}}{\sigma} \left( \frac{S_n}{n} - \mu \right) = \frac{S_n - n\mu}{\sigma\sqrt{n}} \]
е в сила сходимост по разпределение
\[ Z_n \xrightarrow[n \to \infty]{d} Z \sim N(0,1), \]
където $Z$ е стандартно нормално разпределена случайна величина (със средно $0$ и дисперсия $1$).
\end{keythm}

Забележете разликата със закона за големите числа: тук изискваме съществуването на втория момент (дисперсията), което моментално влече съществуването на очакването.

\subsection{Коментари и интуиция зад ЦГТ}

Какво ни казва най-грубо този резултат? От алгебрична гледна точка можем да изразим разликата като:
\[ \frac{S_n}{n} - \mu \approx \frac{\sigma}{\sqrt{n}} Z \]
Това означава, че грешката между емпиричното средно $S_n/n$ и истинското средно $\mu$ е от порядъка на $\frac{\sigma}{\sqrt{n}}$, умножено по някаква случайна пертурбация $Z$, концентрирана около нулата (тъй като разпределението на $Z$ е нормално с връх в нулата). Скоростта на сходимост към нулата се предопределя изцяло от детерминистичното количество $\frac{\sigma}{\sqrt{n}}$, докато флуктуациите идват от $Z$.

\textbf{Универсалност:} Защо този резултат е толкова важен? Защото ние никъде не специфицираме какво е разпределението на изходните случайни величини $X_i$. Те могат да бъдат дискретни, непрекъснати, с най-разнообразни плътности или вероятностни маси. Единственото, което зависи, е да познаваме техните две основни характеристики: средно $\mu$ и дисперсия $\sigma^2$. Веднъж фиксирани ли са те, кумулативният ефект при събиране винаги асимптотично клони към нормално разпределение.

\textbf{Сходимост по разпределение:} Какво означава $Z_n \xrightarrow{d} Z$? Това означава сходимост на функциите на разпределение. Тъй като функцията на разпределение на стандартното нормално разпределение $\Phi(z)$ е непрекъсната за всяко $z \in \mathbb{R}$, имаме:
\[ \mathbb{P}(Z_n < b) \xrightarrow[n \to \infty]{} \mathbb{P}(Z < b) = \Phi(b). \]
По същия начин, за всеки интервал $(a, b)$:
\[ \mathbb{P}(a \le Z_n < b) \xrightarrow[n \to \infty]{} \mathbb{P}(a \le Z < b) = \Phi(b) - \Phi(a). \]
Поради непрекъснатостта на $Z$, няма значение дали интервалите са отворени или затворени.

Ако преработим неравенството $a \le \frac{S_n - n\mu}{\sigma\sqrt{n}} < b$, получаваме:
\[ \mathbb{P}\left( n\mu + a\sigma\sqrt{n} \le S_n < n\mu + b\sigma\sqrt{n} \right) \xrightarrow[n \to \infty]{} \Phi(b) - \Phi(a). \]
Ако вземем много широки граници, например $a = -100$ и $b = 100$, интегралът на нормалната плътност в тези граници е на практика единица, следователно вероятността парциалната сума $S_n$ да попадне в този обхват също е приблизително 1.

\subsection{Примери за приложение на теорема~\ref{thm:clt}}

\begin{example}[Игра на дърпане на въже]\label{ex:11-1}
Връщаме се към \crosslecture{ex:tug-of-war}{модела с дърпане на въже от предходната лекция}. Около въжето се подреждат хора отляво и отдясно. Всеки човек има вероятност $1/2$ да застане отдясно и $1/2$ да застане отляво, и всички дърпат с еднаква сила. Нека $X_j$ бъде приносът на $j$-тия човек, като $X_j = 1$ с вероятност $1/2$ (дърпа надясно) и $X_j = -1$ с вероятност $1/2$ (дърпа наляво). Резултантната сила е:
\[ S_n = \Delta_n - \Lambda_n = \sum_{j=1}^n X_j \]
където $\Delta_n$ и $\Lambda_n$ са съответно броят хора отдясно и отляво. Ясно е, че очакването е $\E[X_1] = 0$. Дисперсията е:
\[ \Var[X_1] = \E[X_1^2] = (1)^2 \frac{1}{2} + (-1)^2 \frac{1}{2} = 1. \]
Следователно $\mu=0$ и $\sigma=1$. По ЦГТ имаме:
\[ \frac{S_n}{\sqrt{n}} \approx Z \sim N(0,1). \]
Каква е вероятността дясната страна да победи (т.е. $S_n > 0$)?
\[ \mathbb{P}(S_n > 0) = \mathbb{P}\left(\frac{S_n}{\sqrt{n}} > 0\right) \approx \mathbb{P}(Z > 0) = \frac{1}{2}. \]
Въпреки че съществува вероятност за равенство ($S_n = 0$), тя клони към нула при $n \to \infty$, тъй като разпределението се апроксимира от непрекъснато такова.

\end{example}

\begin{example}[Непрекъсната случайна величина]\label{ex:11-2}
Нека сега $Y_1$ бъде непрекъсната случайна величина с плътност
\[
f_{Y_1}(y) =
\begin{cases}
\tfrac{1}{2} e^{-y}, & y > 0 \\[2pt]
\tfrac{1}{4}, & y \in (-2, 0) \\[2pt]
0, & \text{иначе.}
\end{cases}
\]
Плътността е подчертано несиметрична — експоненциално отмаляване надясно, константа наляво — и въпреки това $\E[Y_1] = 0$.\footnote{Директно пресмятане дава $\E[Y_1] = \tfrac12 - \tfrac12 = 0$ и $\E[Y_1^2] = \tfrac53$, тоест $\Var[Y_1] = \tfrac53$. (В лекцията дисперсията беше спомената предположително като единица.) Стойността ѝ не влияе на извода: нормирането в ЦГТ се извършва с $\sigma$, каквото и да е то.}
Ако образуваме сумата $S_n' = \sum_{j=1}^n Y_j$, поведението на $\frac{S_n'}{\sqrt{n}}$ ще бъде абсолютно същото като това на дискретните $\pm 1$ стъпки от предния пример. То отново клони към $Z \sim N(0,1)$. В граничния случай изобщо няма разлика между бинарно дискретно разпределение и непрекъснато разпределение, когато става въпрос за суми от голям брой независими екземпляри.

\end{example}

\section{Функции на моментите (Moment Generating Functions)}

За да докажем строго (или поне идейно) Централната гранична теорема~\ref{thm:clt}, ни е необходим мощен математически инструмент — функциите на моментите.

\begin{defn}
Нека $X$ е случайна величина. Функцията на моментите на $X$, означена с $M_X(t)$, се дефинира като математическото очакване на експонентата от $tX$:
\[ M_X(t) = \E[e^{tX}] \]
за всички $t$, за които това очакване съществува (т.е. е крайно), обикновено в някакъв интервал $|t| < r_0$ ($r_0 > 0$). За дискретни случайни величини това се пресмята като:
\[ M_X(t) = \sum_{i} e^{t x_i} \mathbb{P}(X = x_i). \]
За непрекъснати случайни величини се пресмята чрез интеграл върху плътността:
\[ M_X(t) = \int_{-\infty}^{\infty} e^{tx} f_X(x) dx. \]
\end{defn}

\begin{example}[Равномерно разпределение в интервала от 0 до 1]\label{ex:11-3}
Случайната величина $X \sim U(0,1)$ има плътност $f_X(x) = 1$ за $x \in [0,1]$ и 0 извън този интервал. Функцията на моментите е:
\[ M_X(t) = \int_0^1 e^{tx} \cdot 1 \, dx = \left. \frac{e^{tx}}{t} \right|_{0}^{1} = \frac{e^t - 1}{t}. \]
Въпреки че функцията изглежда недефинирана за $t=0$, по правилото на Лопитал или развитие в ред на Тейлър, лесно се вижда, че $\lim_{t \to 0} \frac{e^t - 1}{t} = 1$, което е точно $\E[e^0] = 1$. Следователно $M_X(t)$ е добре дефинирана за всяко реално $t$ (т.е. $r_0 = \infty$).

\end{example}

\begin{example}[Експоненциално разпределение]\label{ex:11-4}
Нека $X \sim \Exponential(1)$ с плътност $f_X(x) = e^{-x}$ за $x \ge 0$.
\[ M_X(t) = \int_0^\infty e^{tx} e^{-x} dx = \int_0^\infty e^{-x(1-t)} dx. \]
Този интеграл съществува и е равен на:
\[ M_X(t) = \frac{1}{1-t}, \quad \text{за } t < 1. \]
Ако $t \ge 1$, експонентата не затихва и интегралът е безкрайност. В този случай функцията на моментите съществува сигурно в интервала $(-1, 1)$, тоест за $|t| < 1$.

\end{example}

\subsection{Моменти на случайна величина}
Преди да преминем към свойствата, нека дефинираме самите моменти, откъдето идва името на функцията:
\begin{itemize}
    \item Обикновен момент от ред $k$: $\E[X^k]$. За $k=1$ това е очакването $\mu$.
    \item Абсолютен момент от ред $k$: $\E[|X|^k]$.
    \item Централен момент от ред $k$: $\E[(X - \mu)^k]$. За $k=2$ това е дисперсията $\Var[X] = \sigma^2$. Центрираме величината, изваждайки нейното средно.
    \item Абсолютен централен момент от ред $k$: $\E[|X - \mu|^k]$.
\end{itemize}
В практиката най-често се използват моментите до четвърти ред (очакване, дисперсия, асиметрия/skewness и ексцес/kurtosis).

\subsection{Свойства на функциите на моментите}

\begin{enumerate}
    \item \textbf{Стойност в нулата:} $M_X(0) = \E[e^0] = 1$.
    \item \textbf{Връзка с моментите:} Използвайки развитието на експоненциалната функция в ред на Тейлър, $e^{tX} = \sum_{k=0}^\infty \frac{t^k X^k}{k!}$, и разменяйки очакването и сумата, получаваме:
    \[ M_X(t) = \E \left[ \sum_{k=0}^\infty \frac{t^k X^k}{k!} \right] = \sum_{k=0}^\infty \frac{t^k}{k!} \E[X^k]. \]
    Това представяне кодира моментите на $X$. $k$-тата производна на $M_X(t)$ в точката $t=0$ ни дава точно $k$-тия момент:
    \[ \left. \frac{d^k}{dt^k} M_X(t) \right|_{t=0} = \E[X^k]. \]
    \item \textbf{Независимост:} Ако $X$ и $Y$ са независими случайни величини, тогава:
    \[ M_{X+Y}(t) = \E[e^{t(X+Y)}] = \E[e^{tX} e^{tY}] = \E[e^{tX}] \E[e^{tY}] = M_X(t) M_Y(t). \]
    Това свойство прави функциите на моментите изключително удобни при работа със суми на независими величини.
    \item \textbf{Линейна трансформация:} Ако $Y = aX + b$, тогава:
    \[ M_Y(t) = \E[e^{t(aX+b)}] = e^{bt} \E[e^{atX}] = e^{bt} M_X(at). \]
    \item \textbf{Единственост:} Ако $M_X(t) = M_Y(t)$ за всички $t$ в някаква околност на нулата ($|t| < r_0$), то $X$ и $Y$ имат едно и също разпределение ($X \stackrel{d}{=} Y$).
    \item \textbf{Сходимост (Теорема на Леви-Крамер):} Ако редица от функции на моментите $M_{X_n}(t)$ се схожда поточково към функция на моментите $M_X(t)$ за всички $t \in (-\varepsilon, \varepsilon)$, то съответните случайни величини се схождат по разпределение: $X_n \xrightarrow{d} X$. Това е в основата на доказателството на ЦГТ.
\end{enumerate}

\begin{remark}[Когато функцията на моментите не съществува]\label{rem:charfun}
Съществуват случайни величини с крайни средно и дисперсия, за които функцията на
моментите не съществува. За тях същата работа върши \emph{характеристичната
функция} $\varphi_X(t) = \E[e^{itX}]$, която е дефинирана винаги, тъй като
$|e^{itX}| = 1$. Идеята е сходна, но апаратът, необходим за нея (комплексен
анализ), излиза извън рамките на този курс, затова тук работим само с $M_X(t)$.
\end{remark}

\subsection{Функция на моментите на нормалното разпределение}

Нека $Z \sim N(0,1)$. Плътността е $f_Z(y) = \frac{1}{\sqrt{2\pi}} e^{-y^2/2}$. Изчисляваме $M_Z(s)$:
\[ M_Z(s) = \frac{1}{\sqrt{2\pi}} \int_{-\infty}^{\infty} e^{sy} e^{-\frac{y^2}{2}} dy = \frac{1}{\sqrt{2\pi}} \int_{-\infty}^{\infty} e^{-\frac{1}{2}(y^2 - 2sy)} dy. \]
Допълваме израза в експонентата до точен квадрат:
\[ y^2 - 2sy = y^2 - 2sy + s^2 - s^2 = (y-s)^2 - s^2. \]
Заместваме обратно в интеграла:
\[ M_Z(s) = \frac{1}{\sqrt{2\pi}} \int_{-\infty}^{\infty} e^{-\frac{1}{2}((y-s)^2 - s^2)} dy = e^{\frac{s^2}{2}} \int_{-\infty}^{\infty} \frac{1}{\sqrt{2\pi}} e^{-\frac{(y-s)^2}{2}} dy. \]
Интегралът отдясно не е нищо друго освен лицето под плътността на нормално разпределение със средно $s$ и дисперсия 1. Тъй като това е валидна плътност, интегралът е точно 1. Следователно:
\[ M_Z(s) = e^{\frac{s^2}{2}}. \]

Сега, ако имаме произволна нормална случайна величина $X \sim N(\mu, \sigma^2)$, ние знаем от свойствата на нормалното разпределение (структурна теорема), че $X$ може да се представи като $X = \mu + \sigma Z$. Прилагайки свойството за линейна трансформация, получаваме:
\[ M_X(t) = e^{\mu t} M_Z(\sigma t) = e^{\mu t} e^{\frac{(\sigma t)^2}{2}} = e^{\mu t + \frac{\sigma^2}{2} t^2}. \]

\section{Доказателство на теорема~\ref{thm:clt}}

\textbf{Цел:} Да докажем, че $Z_n = \frac{S_n - n\mu}{\sigma \sqrt{n}} \xrightarrow[n \to \infty]{d} Z \sim N(0,1)$.

Ще преработим израза за $Z_n$, за да го приведем в по-удобен вид за работа с функции на моментите. Разделяме всяко от събираемите:
\[ Z_n = \frac{1}{\sqrt{n}} \frac{S_n - n\mu}{\sigma} = \frac{1}{\sqrt{n}} \sum_{j=1}^n \frac{X_j - \mu}{\sigma} = \frac{1}{\sqrt{n}} \sum_{j=1}^n Y_j \]
където сме дефинирали $Y_j = \frac{X_j - \mu}{\sigma}$.
Очевидно $Y_j$ са независими и еднакво разпределени случайни величини. За тяхното средно и дисперсия имаме:
\[ \E[Y_1] = \frac{\E[X_1] - \mu}{\sigma} = 0, \quad \Var[Y_1] = \frac{\Var[X_1]}{\sigma^2} = 1. \]
Ще допуснем за улеснение, че функцията на моментите $M_{X_1}(t)$ е добре дефинирана за всички $t$. Следователно, функцията на моментите $M_{Y_1}(t)$ също съществува.

Търсим функцията на моментите на $Z_n$:
\[ M_{Z_n}(t) = M_{\frac{1}{\sqrt{n}} \sum_{j=1}^n Y_j}(t) = M_{\sum_{j=1}^n Y_j}\left(\frac{t}{\sqrt{n}}\right). \]
Тъй като $Y_j$ са независими, функцията на моментите на сумата се разпада в произведение от функциите на моментите на отделните събираеми:
\[ M_{Z_n}(t) = \prod_{j=1}^n M_{Y_j}\left(\frac{t}{\sqrt{n}}\right) = \left( M_{Y_1}\left(\frac{t}{\sqrt{n}}\right) \right)^n. \]
Сега трябва да разберем поведението на $M_{Y_1}\left(\frac{t}{\sqrt{n}}\right)$. Развиваме $M_{Y_1}(s)$ в ред на Тейлър около нулата:
\[ M_{Y_1}(s) = 1 + s\E[Y_1] + \frac{s^2}{2}\E[Y_1^2] + o(s^2). \]
Знаем, че $\E[Y_1] = 0$. За втория момент, тъй като очакването е нула, той съвпада с дисперсията, тоест $\E[Y_1^2] = 1$. Следователно:
\[ M_{Y_1}(s) = 1 + \frac{s^2}{2} + o(s^2). \]
Заместваме $s = \frac{t}{\sqrt{n}}$:
\[ M_{Y_1}\left(\frac{t}{\sqrt{n}}\right) = 1 + \frac{t^2}{2n} + o\left(\frac{t^2}{n}\right). \]
Връщайки се обратно към $M_{Z_n}(t)$, получаваме:
\[ M_{Z_n}(t) = \left( 1 + \frac{t^2 / 2}{n} + o\left(\frac{1}{n}\right) \right)^n. \]
За фиксирано $t$, когато $n \to \infty$, това е добре познатата граница за експоненциалната функция $\lim_{n \to \infty} (1 + \frac{x}{n})^n = e^x$. В нашия случай $x = \frac{t^2}{2}$:
\[ M_{Z_n}(t) \xrightarrow[n \to \infty]{} e^{\frac{t^2}{2}}. \]
Но това е точно функцията на моментите $M_Z(t)$ на стандартното нормално разпределение $Z \sim N(0,1)$. По теоремата за сходимост, тъй като функциите на моментите се схождат, то $Z_n \xrightarrow{d} Z$. Доказателството е завършено.

\section{Приложения и оценка на грешката}

\subsection{ЦГТ за Биномно разпределение}
Нека $X_n \sim \Bin(n, p)$ е биномно разпределена случайна величина. Тя може да бъде представена като сума от независими Бернулиеви случайни величини:
\[ X_n = \sum_{j=1}^n Y_j, \quad Y_j \sim \Ber(p). \]
За Бернулиево разпределение имаме очакване $\mu = p$ и дисперсия $\sigma^2 = p(1-p)$. Следователно, прилагайки ЦГТ за тази сума, получаваме:
\[ \frac{X_n - np}{\sqrt{np(1-p)}} \xrightarrow[n \to \infty]{d} Z. \]
Това ни позволява да оценяваме лесно вероятности от вида $\mathbb{P}(X_n \ge k_n)$, където $k_n = np + a$:
\[ \mathbb{P}(X_n \ge k_n) = \mathbb{P}\left( \frac{X_n - np}{\sqrt{np(1-p)}} \ge \frac{k_n - np}{\sqrt{np(1-p)}} \right) \approx \mathbb{P}\left( Z \ge \frac{k_n - np}{\sqrt{np(1-p)}} \right). \]

\subsection{Оценка на грешката при статистически проучвания и Монте Карло}
Продължаваме \crosslecture{ex:monte-carlo}{приложението на закона за големите числа към Монте Карло от предходната лекция}. Ако интервюираме $N$ човека за кого ще гласуват, или пък използваме метод Монте Карло за намиране на лице на фигура чрез хвърляне на случайни точки, ние реално правим серия от експерименти на Бернули. Истинската вероятност (пропорция) за успех е неизвестно $p$.
Искаме да намерим колко експеримента са нужни, така че емпиричното средно $S_N/N$ да се отклонява от $p$ с повече от $\varepsilon$ с вероятност по-малка от $\delta$. Дефинираме грешката $E_N = \frac{S_N}{N} - p$.
\[ \mathbb{P}(|E_N| > \varepsilon) = \mathbb{P}\left( \frac{|S_N - Np|}{N} > \varepsilon \right) = \mathbb{P}\left( \frac{|S_N - Np|}{\sqrt{N p(1-p)}} > \frac{\varepsilon \sqrt{N}}{\sqrt{p(1-p)}} \right). \]
Тъй като $p$ е неизвестно, ние максимизираме дисперсията. Знаем, че $p(1-p) \le 1/4$ за $p \in (0,1)$ (достига се при $p=1/2$). Оттук $\sqrt{p(1-p)} \le 1/2$. Следователно изразът отдясно става:
\[ \frac{\varepsilon \sqrt{N}}{\sqrt{p(1-p)}} \ge \frac{\varepsilon \sqrt{N}}{1/2} = 2\varepsilon\sqrt{N}. \]
Следователно, оценявайки вероятността отгоре, заместваме с по-малко число в условието на нормалното разпределение:
\[ \mathbb{P}(|E_N| > \varepsilon) \le \mathbb{P}(|Z_N| > 2\varepsilon\sqrt{N}) \approx \mathbb{P}(|Z| > 2\varepsilon\sqrt{N}). \]
Това може да се изрази чрез интеграла от плътността на нормалното разпределение:
\[ \mathbb{P}(|Z| > 2\varepsilon\sqrt{N}) = \frac{1}{\sqrt{2\pi}} \int_{|y| > 2\varepsilon\sqrt{N}} e^{-\frac{y^2}{2}} dy. \]
Ако положим $2\varepsilon\sqrt{N} = A$, тоест $\varepsilon = \frac{A}{2\sqrt{N}}$, можем да запишем:
\[ \mathbb{P}\left(|E_N| > \frac{A}{2\sqrt{N}}\right) \le \frac{1}{\sqrt{2\pi}} \int_{|y|>A} e^{-\frac{y^2}{2}} dy \le \frac{2}{\sqrt{2\pi}} \frac{e^{-\frac{A^2}{2}}}{A} \le \delta. \]
Множителят $2$ идва от това, че интегрираме по двете опашки: $\int_A^\infty e^{-y^2/2}\,dy \le e^{-A^2/2}/A$.
По този начин, ако фиксираме нивото на достоверност $\delta$ (напр. $0{,}01$) и желаната грешка $\varepsilon$, можем да решим горното уравнение, за да изчислим необходимия обем извадка $N$. Това е математическата обосновка зад факта, че социологическите проучвания обикновено включват извадка от около 1000 до 2000 души, за да постигнат надеждни резултати.

Важно е да отбележим колко точно е нашето приближение. Знакът $\approx$ между функцията на разпределение на $Z_n$ и на стандартното нормално разпределение може да бъде строго оценен.

\begin{thm}[Неравенство на Бери\,--\,Есен]\label{thm:berry-esseen}
За всяко $x \in \mathbb{R}$
\[ \left| \mathbb{P}(Z_n < x) - \Phi(x) \right| \le \frac{C \cdot \E[|X_1 - \mu|^3]}{\sigma^3 \sqrt{n}}, \]
където $C$ е абсолютна константа (около $0{,}47$).
\end{thm}
 Това показва, че грешката на самата апроксимация по ЦГТ намалява със скорост $1/\sqrt{n}$ и зависи от третия централен момент на случайните величини.

\section{Функция на моментите на Гама разпределение}

Като приложение, демонстриращо гъвкавостта на този инструмент, ще изчислим функцията на моментите за Гама разпределение. Нека $X \sim \Gamma(\alpha, \beta)$. Плътността е дадена от:
\[ f_X(x) = \frac{\beta^\alpha}{\Gamma(\alpha)} x^{\alpha-1} e^{-\beta x}, \quad \text{за } x > 0. \]
Функцията на моментите се задава с интеграла:
\[ M_X(t) = \int_0^\infty e^{tx} \frac{\beta^\alpha}{\Gamma(\alpha)} x^{\alpha-1} e^{-\beta x} dx = \frac{\beta^\alpha}{\Gamma(\alpha)} \int_0^\infty x^{\alpha-1} e^{-x(\beta-t)} dx. \]
За да решим интеграла, полагаме нова променлива $y = x(\beta-t)$, което предполага, че $t < \beta$ (за да е сходящ интегралът). Тогава $x = \frac{y}{\beta-t}$ и $dx = \frac{dy}{\beta-t}$. Замествайки:
\[ M_X(t) = \frac{\beta^\alpha}{\Gamma(\alpha)} \int_0^\infty \left( \frac{y}{\beta-t} \right)^{\alpha-1} e^{-y} \frac{dy}{\beta-t} = \frac{\beta^\alpha}{\Gamma(\alpha)(\beta-t)^\alpha} \int_0^\infty y^{\alpha-1} e^{-y} dy. \]
Забелязваме, че подинтегралният израз заедно с $1/\Gamma(\alpha)$ формира плътността на Гама разпределение с параметри $\Gamma(\alpha, 1)$:
\[ \frac{1}{\Gamma(\alpha)} y^{\alpha-1} e^{-y} \]
Тъй като интегралът от всяка валидна плътност в граници от $0$ до $\infty$ е 1, то:
\[ \int_0^\infty y^{\alpha-1} e^{-y} dy = \Gamma(\alpha). \]
Тогава изразът се опростява до:
\[ M_X(t) = \frac{\beta^\alpha}{(\beta-t)^\alpha}, \quad \text{за } t < \beta. \]

С помощта на функцията на моментите лесно пресмятаме самите моменти. За да намерим очакването, диференцираме един път:
\[ M_X'(t) = \alpha \beta^\alpha (\beta-t)^{-\alpha-1} \]
Заместваме $t=0$:
\[ \E[X] = M_X'(0) = \alpha \beta^\alpha \beta^{-\alpha-1} = \frac{\alpha}{\beta}. \]
За втория момент намираме втората производна:
\[ M_X''(t) = \alpha(\alpha+1) \beta^\alpha (\beta-t)^{-\alpha-2} \]
При $t=0$:
\[ \E[X^2] = M_X''(0) = \alpha(\alpha+1) \beta^\alpha \beta^{-\alpha-2} = \frac{\alpha(\alpha+1)}{\beta^2}. \]
Тази техника спестява тежкото директно интегриране и е мощен похват в теорията на вероятностите.

\section{Задачи}

\begin{enumerate}
    \item За $Z \sim N(0,1)$ проверете чрез допълване до квадрат, смяна на променливата и свойството, че интегралът от плътността е единица, че
    \[
    M_Z(s)=\frac{1}{\sqrt{2\pi}}\int_{-\infty}^{\infty}e^{sy-y^2/2}\,dy=e^{s^2/2}.
    \]

    \item В модела с дърпане на въже от пример~\ref{ex:11-1} променете условието за победа: дясната страна печели само ако резултантната сила е по-голяма от $100$. Помислете как Централната гранична теорема се прилага към новото събитие.

    \item За $X\sim\Gamma(\alpha,\beta)$ използвайте
    \[
    M_X(t)=\frac{\beta^\alpha}{(\beta-t)^\alpha},\qquad t<\beta,
    \]
    и последователно диференциране в $t=0$, за да пресметнете моментите от по-висок ред.
\end{enumerate}
